%=============================================================
% RESULTS AND DISCUSSION
%=============================================================

\subsection{Reference Front and Coupling}
\label{sec:front}
\label{sec:sensitivity}
Of the $12{,}000$ grid designs, $3{,}150$ break a memory limit and $378$
($3.15\%$) meet every constraint. The latency limit fails $6{,}943$ designs,
the deadline-miss limit $3{,}732$, and the accuracy floor $4{,}200$, with
overlaps among them. The accuracy--energy front of the
acceptable designs has four members, with accuracy from $0.7672$ to $0.8246$
and energy from $9{,}610$ to $14{,}782\,\mu$J per inference. Corner A gives
the front a hypervolume of $752.5$ and corner B $830.7$. The members use
$L=4$ or $5$, weight precision $8$ (three designs) or $16$ (one), and no
clock, priority, or time slice is shared by all four: the clock is
$120$\,MHz for three designs and $180$\,MHz for one, the inference priority is
\emph{below} for two and \emph{equal} for two, and the slices are $10.3$\,ms
and $1.0$\,ms. No design with $b=32$ or $b=4$ is on the front. In this model
the absence of $b=4$ follows from the accuracy loss at $4$ bits and from the
int4 cycle estimate of $1.14$ times the int8 cycles, a literature value that
was not measured (Fig.~\ref{fig:front} shows the front).

\begin{table}[t]
\centering
\caption{Reference Front}
\label{tab:front}
\footnotesize
\resizebox{\columnwidth}{!}{%
\begin{tabular}{|c|c|c|c|l|c|c|c|c|c|}
\hline
\rowcolor{tblhead}
\textbf{$L$} & \textbf{$w$} & \textbf{$b$} & \textbf{$\tau$} & \textbf{$\pi$} & \textbf{$f$} & \textbf{Val.} & \textbf{Test} & \textbf{$E$} & \textbf{$\ell_{95}$} \\
\hline
4 & 0.25 & 8 & 10.3 & equal & 120 & 0.7672 & 0.7735 & 9{,}610 & 56.6 \\
\hline
5 & 0.25 & 8 & 10.3 & equal & 120 & 0.7916 & 0.7830 & 10{,}010 & 69.6 \\
\hline
5 & 0.25 & 16 & 1.0 & below & 120 & 0.7946 & 0.7838 & 10{,}718 & 96.3 \\
\hline
4 & 0.4375 & 8 & 1.0 & below & 180 & 0.8246 & 0.8229 & 14{,}782 & 92.8 \\
\hline
\end{tabular}}

\vspace{2pt}
{\footnotesize $\tau$ in ms, $f$ in MHz, $E$ in $\mu$J per inference, $\ell_{95}$ in ms; accuracy is validation (table) and test.}
\end{table}


\begin{table}[t]
\centering
\caption{Reference Front Under Changed Scenarios}
\label{tab:sensitivity}
\footnotesize
\resizebox{\columnwidth}{!}{%
\begin{tabular}{|l|c|c|c|c|l|}
\hline
\rowcolor{tblhead}
\textbf{Case} & \textbf{Acceptable} & \textbf{Front} & \textbf{$b$} & \textbf{$L$} & \textbf{$f$ (MHz)} \\
\hline
Nominal & 378 & 4 & 16, 8 & 4, 5 & 120, 180 \\
\hline
Latency 50 ms & 143 & 2 & 8 & 4, 5 & 150 \\
\hline
Latency 200 ms & 731 & 5 & 16, 8 & 4, 5 & 60, 120, 150 \\
\hline
Miss limit 0\% & 333 & 4 & 16, 8 & 4, 5 & 120, 180 \\
\hline
Miss limit 5\% & 387 & 4 & 16, 8 & 4, 5 & 120, 180 \\
\hline
Mean latency & 663 & 5 & 16, 8 & 4, 5 & 120, 150 \\
\hline
Stop-mode idle & 378 & 4 & 16, 8 & 4, 5 & 180 \\
\hline
Switch and tick $\times0.5$ & 390 & 4 & 16, 8 & 4, 5 & 120, 180 \\
\hline
Switch and tick $\times2$ & 378 & 4 & 16, 8 & 4, 5 & 120, 180 \\
\hline
Int4 overhead $\times0.5$ & 391 & 5 & 16, 4, 8 & 4, 5 & 120, 180 \\
\hline
Task set, 20\% load & 299 & 3 & 16, 8 & 4, 5 & 120, 150 \\
\hline
Task set, 50\% load & 99 & 2 & 8 & 4, 5 & 180 \\
\hline
\end{tabular}}
\end{table}


Sensitivity runs recompute the front under $19$ cases, $12$ of them in
the table: latency limits of
$50$, $100$, and $200$\,ms against miss limits of $0$, $1$, and $5\%$, mean in
place of $95$th-percentile latency, Stop-mode idle current, switch and tick
costs scaled by $0.5$ and $2$, the int16 and int4 overheads scaled by $0.5$
and $2$, and two generated task sets at $20$ and $50\%$ periodic utilization
\cite{kim2021harmonic}. The front has $2$ to $6$ designs in every case, always
with $L\in\{4,5\}$ and never with $b=32$. Int4 designs enter only when the
int4 overhead is halved. Tightening the latency limit to $50$\,ms pushes the
front to $b=8$ alone and shrinks the acceptable set to $125$ to $143$
designs, and a $200$\,ms limit admits $60$\,MHz designs and up to $812$
acceptable ones. With Stop-mode idle current, the front moves to
$180$\,MHz alone, the race-to-idle result. The task set at $50\%$
utilization leaves $99$ acceptable designs. The coupling between
scheduler and network therefore shows in which clock and precision survive
each constraint set, and it is stable in the layer count.

\subsection{Method Comparison}
\label{sec:comparison}
\begin{table*}[t]
\centering
\caption{Method Comparison}
\label{tab:methods}
\small

\begin{tabular}{|l|c|c|c|c|c|c|}
\hline
\rowcolor{tblhead}
\textbf{Method} & \textbf{HV$_A$ (\%)} & \textbf{HV$_B$ (\%)} & \textbf{Evals. to 90\%} & \textbf{IGD} & \textbf{$p$} & \textbf{$\delta$} \\
\hline
qNEHVI & 100.0 (100.0--100.0) & 100.0 & 28 & 0 & -- & -- \\
\hline
qNParEGO & 99.6 (96.4--100.0) & 99.7 & 40 & $1.288\times10^{-2}$ & $1.79\times10^{-3}$ & 0.75 \\
\hline
Sequential & 91.2 (91.2--91.2) & 70.3 & 65 & 0.150 & $9.79\times10^{-5}$ & 0.96 \\
\hline
MOTPE & 91.0 (88.7--91.0) & 70.2 & 54 & 0.151 & $1.17\times10^{-4}$ & 0.99 \\
\hline
Random & 74.8 (65.0--84.5) & 58.9 & never & 0.249 & $4.77\times10^{-6}$ & 0.99 \\
\hline
NSGA-II & 56.0 (0.0--74.1) & 28.6 & never & 0.338 & $1.17\times10^{-4}$ & 0.99 \\
\hline
\end{tabular}

\vspace{2pt}
{\footnotesize Median (IQR) over 20 seeds at 100 evaluations. HV is the percentage of the reference-front hypervolume at corner A, and at corner B in the next column. $p$ is the one-sided Wilcoxon test of qNEHVI against the method with Holm correction, and $\delta$ is Cliff's delta.}
\end{table*}


At $100$ evaluations, qNEHVI reaches a median of $100\%$ of the reference
hypervolume at corner A, and $18$ of $20$ seeds end at exactly that value;
one ends at $91.2\%$ and another at $99.95\%$ (Fig.~\ref{fig:convergence}). qNParEGO reaches
$99.6\%$ (interquartile range $96.4$ to $100$), MOTPE $91.0\%$, the
two-stage Sequential search $91.2\%$, Random $74.8\%$, and NSGA-II $56.0\%$.
Corner B gives the same order, $100$, $99.7$, $70.2$, $70.3$, $58.9$, and
$28.6\%$. The Friedman test rejects equal methods at both corners ($p=1.2\times10^{-16}$
and $2.6\times10^{-15}$). qNEHVI is also the fastest to $90\%$, at a median of $28$ evaluations
against $40$ for qNParEGO, $54$ for MOTPE, and $65$ for Sequential; Random and
NSGA-II reach it in only one seed of $20$. Against Random, NSGA-II, MOTPE,
and Sequential, the Holm-corrected one-sided $p$-values lie between
$4.8\times10^{-6}$ and $1.2\times10^{-4}$, where the smallest value $20$ pairs can reach is
$9.5\times10^{-7}$, and Cliff's $\delta$ is $0.96$ to $0.99$. Against
qNParEGO, $p=1.8\times10^{-3}$ and $\delta=0.75$ at corner A, with a median
difference of $0.2$ points: the two joint Bayesian methods are close, and the
gap, although consistent, is small.

\begin{figure}[t]
\centering
\includegraphics[width=\columnwidth]{fig_convergence.pdf}
\caption{Hypervolume as a percentage of the reference front against
evaluations, median over $20$ seeds with the interquartile range shaded.
The dotted line marks $90\%$.}
\label{fig:convergence}
\end{figure}

The share of evaluations that were not acceptable is $56\%$ for qNEHVI,
$58\%$ for MOTPE, $57\%$ for Sequential, $66\%$ for qNParEGO, $79\%$ for
NSGA-II, and $97\%$ for Random. With $378$ acceptable designs among
$12{,}000$ and a front of four, the space is small for a model-based search, and the margin over
Random is a statement about this problem; a larger space or a longer front
could change it. NSGA-II ran with pymoo defaults and a population of $10$,
and on some seeds it collapses to one architecture; its rank reflects that
configuration and was not tuned.

Proposal time separates the methods in wall-clock terms. The median total
time spent proposing designs over $100$ evaluations is $23$ minutes for
qNEHVI, $33$ for qNParEGO, and $58$ for Sequential, which is $14$, $20$, and $35$
seconds per evaluation, on a fanless laptop. NSGA-II, MOTPE, and Random need
under one second per run. One evaluation of a design costs milliseconds, so
in this problem the Bayesian methods are slower overall; they pay off only
where an evaluation costs more than their proposal time, as a training run
or a board measurement would. In evaluations, qNEHVI reaches $90\%$ after
$28$, which is $0.23\%$ of the grid, and ends at $100$ evaluations, which is
$0.83\%$.

\begin{figure}[t]
\centering
\includegraphics[width=\columnwidth]{fig_front.pdf}
\caption{The reference front and the front found by the median seed of each
method. Markers are offset sideways by a few $\mu$J so that coincident
designs stay visible.}
\label{fig:front}
\end{figure}

\subsection{Joint Against Sequential Search}
\label{sec:sequential}
Sequential stalls at $91.2\%$ at corner A and $70.3\%$ at corner B in all $20$
seeds, so its interquartile range is zero. The cap is structural: its first stage fixes
$180$\,MHz, a $10$\,ms slice, and equal priority, and three of the four
front designs run at $120$\,MHz. No second-stage search can reach them,
because the network was chosen under the wrong clock. Joint qNEHVI exceeds
it by $8.8$ points at the median ($p=9.8\times10^{-5}$, $\delta=0.96$). The
size of the cap depends on the values stage one fixes, and a different fixed
clock would move it; the result shows that fixing the scheduler first can
exclude the best designs, and not by how much in general.

\subsection{Hardware Validation and Test Accuracy}
\label{sec:hardware}
\begin{table}[t]
\centering
\caption{Hardware Validation Error (\%)}
\label{tab:hardware}
\begin{tabular}{|l|c|c|c|c|}
\hline
\rowcolor{tblhead}
\textbf{Models} & \textbf{n} & \textbf{Cycles} & \textbf{Flash} & \textbf{RAM} \\
\hline
Held out & 8 & 4.76 (7.60) & 1.62 (6.30) & 6.64 (8.55) \\
\hline
Calibration & 24 & 2.39 (4.66) & 1.11 (9.46) & 6.83 (14.32) \\
\hline
\end{tabular}

\vspace{2pt}
{\footnotesize Mean absolute error in percent (maximum in brackets) of predicted against measured, STM32F469I-DISCO at 180 MHz.}
\end{table}


\begin{figure}[t]
\centering
\includegraphics[width=\columnwidth]{fig_cycles.pdf}
\caption{Predicted against measured inference cycles on the STM32F469I-DISCO
at $180$\,MHz, for the $24$ calibration benchmarks and the eight held-out
models.}
\label{fig:cycles}
\end{figure}

Eight int8 networks that were not used in calibration ($L=2$ to $5$ at
$w=0.4375$ and $0.8125$) were benchmarked on the same board. Their cycles are
predicted with a mean absolute error of $4.8\%$ and at most $7.6\%$, against
$2.4\%$ and $4.7\%$ on the $24$ calibration benchmarks. The error has a sign:
$7$ of $8$ predictions are low, with a mean of $-4.5\%$ (the calibration set
has $-0.1\%$), so predicted cycles run a few percent below the board. Flash
errs by $1.6\%$ on average and $6.3\%$ at most, and RAM by $6.6\%$ and $8.6\%$.
The four front designs were also evaluated once on the CIFAR-10 test split
of $10{,}000$ images. Test accuracy differs from the table's validation
accuracy by $+0.63$, $-0.17$, $-0.86$, and $-1.08$ points, a mean of $-0.37$. A negative mean is expected when the checkpoint is chosen
on validation accuracy. For comparison, the seed standard deviation is between
$0.44$ and $0.47$ points at $8$, $16$, and $32$ bits.
The int16 and int4 cycle ratios, and the memory of every int16 and int4
design, were not run on silicon, which includes the one front design at
$b=16$.

\subsection{Limitations}
\label{sec:limitations}
Energy comes from datasheet currents and was not measured, and the run
currents are typical values for benchmark code. Cycles and memory were
measured on an F469 board and applied to an F446 target of the same core family.
Table accuracy is one training seed per architecture, with a seed noise of
$0.44$ to $1.26$ points that is larger at int4. The workload is a
synthetic task set with an assumed tick cost, and the front lies on a grid
of $12{,}000$ designs that holds ten values of $\tau$, whereas the searches
treat it as continuous. The two Ax methods receive corner A as objective thresholds and the other
methods use no reference point, which may favor the Ax methods; corner B
gives the same order but uses the same thresholds. Twenty seeds bound each
$p$-value from below at about $10^{-6}$.
